\scanpage{11}
% 1–4
\uline{NB} \deleted{Le plus petit $A_*$} \inserted{a) b) c) impliquent} que $A_*$ est un sous-anneau cosimplicial de $S_*$. Le plus petit $A_*$ satisfaisant à a) b) c) est formé des

% 5–8
\[
\begin{aligned}
 A_n&=\mathrm{Im}\bigl(\mathbb Z[t_1,\ldots,t_n]\hookrightarrow k[[t_1,\ldots,t_n]]\bigr)\\
 &\simeq\Lambda[t_1,\ldots,t_n]
\end{aligned}
\]
où
\[
 \Lambda=\mathrm{Im}(\mathbb Z\unsure{-|\to} k)\quad
 \begin{cases}\simeq\mathbb Z&\text{ou}\\
 \simeq\mathbb Z/n\mathbb Z&\text{\deleted{pour} } n\in\mathbb N^*\ \text{\inserted{conven.}}\end{cases}
\]

% 9–13
Avec les conditions a) b) c) on construit la catégorie \deleted{$\unsure{\mathrm{Sfl}|\mathrm{Sf}}$} $\mathrm{Sfl}_{A_*}$ des \deleted{$A_*$}schémas \inserted{$A_*$-}formels (lisses). Dans cette catégorie, il y a des produits finis ($X_n\simeq(X_1)^n$), et des produits fibrés dans le cas ``transversal''.

% 13, left
\margin{Le foncteur ``oubli de $A_*$'' commute à ces lim $\leftarrow$}

% 14–18
\inserted{Soit $k_0=A_0\subset S_0=k$.} On obtient un foncteur $W_{A_*}$ de la cat. des $k_0$-modules libres de r.f. (et $=$ projectifs de type fini ?) vers les $A_*$-schémas formels, par
\[
 M_0\mapsto\mathrm{Spf}\,\widehat{\mathrm{Sym}}{}_k^*(M_{0k})
\]
(où $M_{0k}=M_0\otimes_{k_0}k$). C'est $=$ un foncteur à valeurs dans les objets groupes de $\mathrm{Sfl}_{A_*}$ i.e. vers les groupes $A_*$-formels.

% 14–18, left upper
% slanted
\margin{\deleted{\unsure{Notion} plus gén. $A_*$-schéma (fibré) sur $A_*\unsure{X|Y}$ et section sur $\unsure{A_X|A_Y}$, on a (suite \unsure{exacte|générale}) .\illegible{1}}}

% 14–18, left lower
% slanted
\margin{Prendre plus gén. des modules proj. de t. fini sur un $O_{A_*}(X)$ ..}

% 17–24, outer edge
% margin note written sideways
\margin{$O_{A_*}(X)\subset O(X)$. $\mathrm{Hom}_{A_*\text{ form}}(X,X_1)\subset\mathrm{Hom}_{k\text{-form}}(X,X_1)$ Modules \inserted{(lisses de t.f.)} $A_*$-formels sur les schémas $A_*$-formels .... \deleted{\unsure{Le|Le Sp}}\inserted{fibré} vectoriel $A_*$ formel associé $\overline V_{A_*}(M)$, ou $W_{A_*}(M)$}

% 19–25
Plus gén., si $\unsure{\mathrm{Lisc}|\mathrm{Linc}|\mathrm{Lin}}(k_0)$ est la catégorie des $k_0$-schémas augmentés \inserted{$k_0$-}isomorphes à \deleted{\unsure{Spec|Sp}} des \deleted{\unsure{E|k}} $\mathrm{Sp}\ k_0[\unsure{\mathrm{Sym}_k^*(M_{k_0})|\mathrm{Sym}^*(M_{k_0})|\widehat{\mathrm{Sym}}_k^*(M_0)}]$, \deleted{\illegible{1}}, il y a un foncteur $X\unsure{-|\to}\widehat X^{(A_*)}$ de $\mathrm{Lisc}(k_0)$ vers les $A_*$-schémas formels, commutant aux produits finis, et transformant objets groupes en objets groupes ....
% archive sheet number: 10
